\chapter[Background and Literature Review]{Background and Literature Review}
\label{chapter:background}

\begin{quote}\itshape
The objective of this chapter is to establish the background and literature
context for this thesis. The chapter begins by reviewing embodied AI and
robotics. It then examines LLM agents and robot software and draws together the
literature relevant to robot-specific software synthesis and validation.
\end{quote}

Embodied and generalist robot intelligence seeks not only to interpret a task,
but also to transform that interpretation into actions that a particular robot
can execute. Although recent generalist models have expanded the range of
tasks, environments, and embodiments addressed through shared representations
and policies, physical execution remains conditioned by each robot's
morphology, action space, control hierarchy, and software interface
\citep{brooks1991intelligence,brohan2022rt1,openxembodiment2024}. The same
task-level command may therefore require substantially different
implementations across robotic platforms. Consequently, generality at the
model level does not, by itself, establish generality at deployment: a
practical mechanism is still required to connect general task intent to
reliable robot-specific execution.

This chapter reviews the relevant literature through two linked questions.
First, within which physical and software structures does robot intelligence
operate? Second, can a large language model (LLM) agent generate the software
required to connect general task intent with executable behaviour on a
particular robot? Three distinctions are maintained throughout the review.
Invoking a supplied capability interface is distinguished from synthesising
either a new capability or the robot-specific driver that implements it. A
capability is treated as a bounded and reusable unit of functionality, rather
than as a complete task-level skill. Finally, successful execution in
simulation is interpreted as evidence about the simulated system and not as
evidence of hardware performance. Section~\ref{sec:embodied-ai-robotics} establishes the embodied,
control, interface, and morphology context, whereas Section~\ref{sec:llm-agents-robot-software} examines
LLM-based robot programming, software synthesis, and behavioural validation.

\section{Embodied AI and Robotics}
\label{sec:embodied-ai-robotics}

\subsection{Embodiment and Action}
\label{sec:embodiment-action}

Embodiment is used in the literature to describe several relationships between
an agent, its body, and its environment, rather than the mere possession of a
physical or simulated body \citep{ziemke2003embodiment}.
\citet{brooks1991intelligence} challenged architectures that treated
intelligence primarily as centralised symbolic representation and instead
located autonomous behaviour in situated perception and action. Later reviews
describe embodied cognition as activity shaped by an agent's bodily capacities
and practical interaction with its surroundings
\citep{anderson2003fieldguide}. This thesis therefore uses embodiment in an
engineering sense: intelligent behaviour depends on the coupling among a
robot's body, sensorimotor apparatus, controller, and environment.

The central mechanism is sensorimotor coupling. When a robot acts, it changes
its own state or the surrounding environment, which in turn changes the sensory
information available for the next decision. Dynamical accounts accordingly
analyse behaviour as a trajectory of the coupled agent--environment system
rather than as the output of an isolated controller
\citep{beer1995dynamical}. Robotics studies likewise show that active movement
can organise sensory information and simplify the problem presented to the
controller \citep{pfeifer1997sensorymotor}. More recent reviews preserve this
controller--body--environment perspective while distinguishing interactions
that occur across different physical and computational timescales
\citep{iida2023timescales}. This thesis adopts that mechanism without taking a
position on whether internal representations are necessary for intelligence.

Action must nevertheless be specified at an abstraction level before it can be
executed. In robotics, an action may refer to a desired effect, an end-effector
reference, a joint-space target, an actuator command, or a temporally extended
movement representation; these objects are related but not interchangeable
\citep{zech2019action}. Operational-space control, for example, formulates
end-effector motion and force separately from the underlying joint dynamics
\citep{khatib1987unified}, whereas dynamical movement primitives encode
adaptable goal-directed trajectories \citep{ijspeert2013dmp}. Learning-based
approaches often express a policy as a mapping from observed state to action,
but the meaning of the action remains determined by the selected representation
and robot context \citep{argall2009survey,kober2013reinforcement}.

This range of meanings makes it necessary to distinguish an intended action,
an issued command, actuator-driven execution, and the resulting effect.
\citet{zech2019action} treat anticipated effects and mechanical effectiveness
as central elements of robotic action, while also noting that no single
definition is accepted across the field. This thesis does not impose such a
definition universally; instead, it uses the narrower empirical distinction
that a command is not evidence of its intended outcome until the relevant
post-action state is observed on the declared execution route. Because command
semantics depend on a robot configuration---including its joint ordering,
coordinate conventions, actuation, limits, and feedback signals---the same
task-level intention may require different executable mappings across robots.
The next subsection examines how control hierarchies organise these mappings
across planning, motion generation, and feedback control.

\subsection{Control Hierarchies}
\label{sec:control-hierarchies}

Across reviews of robot architectures, motor control, and hierarchical
reinforcement learning, hierarchical control is not tied to a fixed number of
software layers. It is instead characterised along three recurring dimensions:
functional abstraction, temporal horizon or update rate, and the delegation of
control and information. Classical architecture surveys describe abstract
tasks being progressively refined through planning, executive, and functional
levels, whereas motor-control and learning reviews emphasise partial autonomy,
subtask decomposition, and temporal abstraction
\citep{medeiros1998architectures,merel2019hierarchical,
pateria2021hrlsurvey,acs2026longhorizon}. This thesis therefore defines a
control hierarchy as interacting functional levels in which higher levels
select longer-horizon goals or operations, lower levels realise them at faster
rates, and execution feedback can revise higher-level decisions.

High-level control answers what should be done and in what order. It includes
deliberative planning and monitoring, integrated task-and-motion planning,
behaviour-tree execution, hierarchical policies, and language-conditioned
capability selection
\citep{ingrand2017deliberation,garrett2021tamp,iovino2022btsurvey,
ahn2022saycan,liang2023codeaspolicies}. Low-level control answers how a
requested effect is physically realised through kinematic and collision-aware
planning, trajectory generation, and operational-, joint-, or actuator-space
feedback control
\citep{khatib1987unified,lynch2017modern,ijspeert2013dmp}. Although end-to-end
vision--language--action policies may combine several of these mappings
\citep{brohan2023rt2,black2024pi0}, recent hierarchical VLA systems retain
explicit semantic subgoal selection above low-level action generation
\citep{shi2025hirobot,hu2026orchestrating}.

The robot capability interface sits at the delegation boundary between these
functions. AutoAdapter-Bench examines the production of this software layer
separately from its use by a high-level controller. The following subsections
establish the planning and interface concepts needed to distinguish these two
roles.

\subsection{Generalist Robot Intelligence}
\label{sec:generalist-robot-intelligence}

Within a hierarchical robot system, generalist intelligence at the high level
concerns interpreting varied task instructions and coordinating reusable
capabilities, rather than directly predicting actuator commands. The methods
below are compared by what the LLM outputs, how that output is constrained,
and what feedback returns after execution.

Direct planners first generated sequences over fixed actions. Zero-Shot
Planner maps free-form steps to admissible actions, while SayCan combines
task-relevance and skill-feasibility scores. Grounded Decoding adds affordance
or safety objectives during token generation, whereas Text2Motion checks
geometric feasibility across an entire skill sequence
\citep{huang2022zeroshotplanner,ahn2022saycan,
huang2023groundeddecoding,lin2023text2motion}. These methods improve
grounding, but still assume supplied skills, policies, and feasibility models.

Closed-loop methods revise plans from execution evidence. Inner Monologue
converts scene, success--failure, and human feedback into text. Interactive
Task Planning separates a planner from a function-calling executor. SayPlan
searches a hierarchical 3D scene graph, validates plans in a scene-graph
simulator, and delegates navigation to a classical planner. DoReMi continuously
monitors language-defined constraints and triggers recovery when plan and
execution diverge
\citep{huang2023innermonologue,li2025interactiveplanning,
rana2023sayplan,guo2024doremi}. Their behaviour therefore depends on feedback
projection, tool contracts, monitoring, and memory---not only on the backbone.

Structured programs make control flow explicit. ProgPrompt organises supplied
action calls into programs with assertions, while LLM-as-BT-Planner generates
behaviour trees and Code-BT converts constrained code into them
\citep{singh2023progprompt,ao2025llmbtplanner,zhang2025codebt}.
These approaches organise an existing action or API vocabulary.
LLM+P and AutoTAMP instead translate language into formal planning inputs for
external solvers \citep{liu2023llmplusp,chen2024autotamp}.
Their results depend on the supplied planning models and the translated task
specification. HBTP uses LLM reasoning to guide behaviour-tree search, while
UHBTP replaces its heuristic with a classical, domain-independent alternative
\citep{cai2025hbtp,chen2025btpg}.

Across these families, the LLM acts as a generator, scorer, feedback
controller, translator, or heuristic. Yet all begin with executable operations
or formal action models. Auto-Adapter targets this preceding boundary by
constructing a capability interface and synthesising its robot-specific
driver. AutoAdapter-Bench separately evaluates driver synthesis and use of a fixed
interface. Its matched tool-use comparisons hold Recursive Context-Aware
Reasoning and Planning (ReCAP) fixed as the high-level controller
\citep{zhang2025recap}. The next section therefore examines interfaces and
drivers.

\subsection{Interfaces and Drivers}
\label{sec:interfaces-drivers}

A high-level controller needs a stable boundary for delegating a task decision
without knowing how a particular robot will realise it. Research on action
representations distinguishes the intended effect of an action from the
end-effector, joint, or actuator commands used to produce that effect
\citep{zech2019action}. Reviews of capability-based frameworks identify a
similar separation, despite varied uses of task, skill, and primitive: an
operation is exposed in terms meaningful to a caller and then grounded in
hardware-facing control \citep{pantano2022capabilitysurvey}. This thesis calls
the caller-facing contract the robot capability interface. It defines what an
operation means, when it may be requested, and what evidence would demonstrate
its outcome. It also fixes conventions such as units and reference frames,
which would otherwise make a call ambiguous, but does not itself execute the
operation.

Execution instead depends on software that connects task-level requests to
the control facilities of one robot, and current systems place this connection
in different artefact types. Within ROS~2, \texttt{ros2\_control} separates
controllers from dynamically loaded hardware components, while SkiROS2 composes
executable skill implementations above ROS
\citep{macenski2022ros2,ros2control2026hardware,mayr2023skiros2}.
LeRobot takes a more compact route by defining a Python \texttt{Robot} class
that exchanges actions and observations through a motor bus or manufacturer
interface \citep{lerobot2026hardware}. These systems expose different software
boundaries, but each connects an operation understood by higher-level software
to a robot-specific means of execution.

The Model Hardware Standard (MHS) proposes a common device interface for AI
agents, developed initially by Anthropic and the HHMI Janelia Research Campus.
Its drivers expose device operations and descriptions through standard access
mechanisms. MHS addresses how equipment is made discoverable and callable;
Auto-Adapter addresses the generation of robot-specific capability
implementations. These roles may be complementary. The cited research preview
motivates standardised integration, but does not provide a directly comparable
evaluation of driver synthesis
\citep{anthropic2026mhs}.

Auto-Adapter places this connection in one generated Python robot-specific
driver for each robot configuration. Its public methods implement the
capability interface, while the code behind them maps requests and feedback to
the representations used by the robot. The driver may invoke existing control
facilities or supply the required motion generation and feedback logic itself
\citep{lynch2017modern}. One artefact may therefore span responsibilities that
other systems divide among several components, while the task-level sequence
remains with the high-level controller.

The division of work between the high-level controller and the driver changes
with capability granularity. A coarse-grained operation leaves fewer decisions
to the high-level controller, but requires the driver to realise more movement
and recovery behind one call. A fine-grained operation permits more
recomposition, but asks the high-level controller to coordinate more steps and
respond to intermediate feedback; neither allocation is uniformly preferable
\citep{zech2019action,pantano2022capabilitysurvey}.

AutoAdapter-Bench compares driver synthesis under a fixed
capability interface and common behavioural tests for each robot. Models
generate and revise the implementation against these shared requirements. The
separate driver-use experiment compares task performance through the same
validated driver, with the ReCAP planning procedure and task conditions held
constant across models.

\section{LLM Agents and Robot Software}
\label{sec:llm-agents-robot-software}

\subsection{LLM Robot Programming}
\label{sec:llm-robot-programming}

LLM robot programming refers here to using an LLM to generate a program that
organises robot behaviour. A separate system executes the program, which
distinguishes this approach from policies that predict continuous actions
directly from observations. Related planning methods ground language
instructions in supplied robot operations. SayCan selects among pre-trained skills, whereas
ProgPrompt completes Python-like task functions using supplied actions and
assertions \citep{ahn2022saycan,singh2023progprompt}. Both rely on
pre-implemented robot operations.

Code as Policies extends code generation from task sequencing to policy logic.
Its generated Python can contain branches, loops, and helper functions. The
program can also process perceptual outputs and parameterise supplied control
APIs. VoxPoser similarly generates code that constructs three-dimensional
value maps, which a motion planner uses to synthesise trajectories
\citep{liang2023codeaspolicies,huang2023voxposer}. Together, these studies show
that code can connect language instructions to robot execution. However, both
approaches still rely on system-provided perception and motion mechanisms.

RoboCodeX broadens this approach through a tree-structured multimodal
code-generation framework. It decomposes high-level instructions into
object-centred manipulation units that include affordances and safety
constraints. Its evaluation includes simulated and physical robot tasks
\citep{mu2024robocodex}. Nevertheless, the evaluated output concerns behaviour
for a requested task. This evidence does not by itself establish that an LLM
can implement a fixed, reusable capability interface for a new robot. The
distinction concerns the generated artefact, rather than a shortcoming of the
method.

This distinction defines the object studied in this thesis. Calling an
existing \texttt{move\_to()} method is capability-interface use, whereas
writing its implementation for a particular robot is driver synthesis.
Combining several capabilities to complete a task is task programming. Prior
work therefore supports executable code as a representation of robot
behaviour, but it commonly begins from available operations. Robot-specific
driver synthesis must additionally understand the surrounding software and
implement the required interface. It must also correct failures found during
execution. The next subsection examines coding agents as a method for
performing this wider software task.

\subsection{Coding Agents}
\label{sec:coding-agents}

Early LLM code generation treated programming as function completion. HumanEval
tests functions generated from docstrings \citep{chen2021codex}. It provides
little context about an existing software system. Success therefore does not
demonstrate repository-level coding.

SWE-bench instead evaluates codebase issues \citep{jimenez2024swebench}.
SWE-agent exposes code and execution results through an agent--computer
interface \citep{yang2024sweagent}. Interaction therefore becomes part of the
coding method.

ReAct alternates reasoning and action so observations can revise later steps
\citep{yao2023react}. CodeAct makes executable code the agent's action
\citep{wang2024codeact}. Auto-Adapter adopts a bounded version of this process.
The agent studies public robot material before producing a driver.
Public simulation feedback can reveal behavioural
faults and guide driver repair. Acceptance requires evaluation against fixed
criteria using observations outside the candidate's control.

Program sketching uses a partial program to express an implementation strategy
while synthesis fills the missing parts \citep{solarlezama2013sketching}. A
controller skeleton similarly supplies structure and control knowledge for
LLM-based driver synthesis. AutoAdapter-Bench provides a robot-specific skeleton
across the models evaluated on that robot. This supplied implementation must be
considered when interpreting model performance.

Feedback has two roles. Reflexion guides later attempts through verbal
reflection, whereas ExpeL extracts reusable knowledge across tasks
\citep{shinn2023reflexion,zhao2024expel}. Within Auto-Adapter, validation
diagnostics guide revisions to the current driver while its capability
definitions and acceptance criteria remain unchanged. This use of feedback
concerns repair within a run.

\subsection{Tool and Driver Synthesis}
\label{sec:tool-driver-synthesis}

This subsection reviews how LLMs synthesise tools for their own use and how
tool synthesis has been applied to robot software. In Auto-Adapter, the
resulting tool takes the form of a robot driver. Language to Rewards uses an
LLM to specify reward parameters and optimises actions online with MuJoCo
model-predictive control \citep{yu2023language2reward}. Eureka instead improves
reward code through simulation feedback and uses reinforcement learning to
train control policies \citep{ma2024eureka}.

LLMCap2Skill generates skills from capability descriptions, user-written
behavioural specifications and retrieved ROS~2 interfaces. Through pySkillUp,
it produces skill interfaces and state machines. Its Gazebo evaluation examines
execution and agreement with the specification, including required manual
corrections. It therefore combines supplied specifications with generated
implementations \citep{vieira2025cap2skill}.

Tripathi et al.\ generate navigation controllers from task specifications and
environment information. Structured tests provide diagnostics for revising
prompts or controller code. Their Webots evaluation addresses tested navigation
behaviour \citep{tripathi2026testdriven}.

CaP-X evaluates manipulation programs across primitive abstraction levels and
feedback conditions, while CaP-Agent0 extracts reusable functions from
successful executions. These programs compose supplied perception and control
services \citep{fu2026capx}.

Tool reuse is also studied outside robotics. LATM separates
tool creation from tool use. It also stores generated Python functions for later
tasks \citep{cai2023latm}. CREATOR first produces a tool description and
implementation, then applies the tool to a target problem
\citep{qian2023creator}. Voyager revises programs after execution failures and
retains successful code in a reusable skill library \citep{wang2023voyager}.

These studies support software reuse and revision through execution feedback,
but differ in the software supplied to the generator and the artefact being
evaluated. A generated reward, task program, reusable function and robot driver
impose different implementation requirements. Comparisons must therefore
account for the supplied control software and behavioural specification, as
well as the observations and criteria used to judge execution.

\subsection{Oracles and Validation}
\label{sec:oracles-validation}

A test oracle determines whether observed behaviour satisfies an expected
outcome. Robot assurance methods differ in how that expectation is specified
and what evidence is available to check it. ROSRV monitors developer-specified
properties over ROS messages and can intercept and modify commands to enforce
them \citep{huang2014rosrv}. HAROS checks user-specified properties against
models extracted from source code and deployment configurations; its analysis
depends on the extracted model and assumptions about node behaviour
\citep{santos2021haros}. SOTER instead monitors system state and switches to a
certified controller when required by its runtime-assurance logic. Its safety
guarantee depends on satisfying the framework's stated conditions
\citep{desai2019soter}. These approaches provide different forms of assurance,
bounded by the properties, models and observations they use.

Robot benchmarks make the evaluation target explicit at the task level.
RoboEval checks generated programs against temporal conditions using traces
from a symbolic simulator across multiple initial states
\citep{hu2024codebotler}. This assesses program logic within the symbolic
model, without establishing physical execution accuracy. ROBEL evaluates
trajectories in simulation or on hardware. Its D'Claw pose task uses a
time-averaged absolute joint error below \(10^\circ\), while separate safety
measures record position, velocity and current violations
\citep{ahn2020robel}. LLMCap2Skill evaluates generated skills in Gazebo against
user-written behavioural specifications, including the manual changes required
to obtain the intended behaviour \citep{vieira2025cap2skill}. These evaluations
differ in both the behaviour observed and the role of human judgement.

For generated robot software, the evaluator's independence is a further
consideration. The candidate implementation should not determine the
measurements or acceptance criteria used to judge it. Diagnostic results can
still guide repair while those criteria remain unchanged. Independence alone
cannot establish that a specification is appropriate or that its tests cover
the intended capability. A passing result therefore supports only the behaviour
examined under the stated conditions; simulation results do not establish
hardware performance or safety.
